\par\textbf{Procedimiento completo para comparar tu solución.} La ley lineal del perfil usa el ángulo relativo al de sustentación nula en radianes. La fuerza y el par se obtienen con sus referencias de cuerda. Para el centro de presión, igualar el momento de la resultante al par respecto al centro aerodinámico.

\begin{enumerate}

\item Convertir el intervalo angular.

\[\fitmath{\mathrm{dalph}=\frac{\left(4--2\right)\cdot \pi}{180}\approx 0.1047198\,\mathrm{rad}}\]

\item Presión dinámica.

\[\fitmath{q=\frac{1.05\cdot {45}^{2}}{2}\approx 1063.125\,\mathrm{Pa}}\]

\item Ley de sustentación del perfil.

\[\fitmath{\mathrm{cl}=2\cdot \pi\cdot \mathrm{dalph}}\]

\[\fitmath{\mathrm{cl}=2\cdot \pi\cdot (0.1047197551)\approx 0.6579736\,}\]

\item Sustentación por envergadura.

\[\fitmath{\mathrm{lift}=q\cdot 1.4\cdot \mathrm{cl}}\]

\[\fitmath{\mathrm{lift}=(1063.125)\cdot 1.4\cdot (0.6579736267)\approx 979.3115\,\mathrm{N/m}}\]

\item Momento por envergadura, negativo de picado.

\[\fitmath{\mathrm{moment}=q\cdot {1.4}^{2}\cdot -0.06}\]

\[\fitmath{\mathrm{moment}=(1063.125)\cdot {1.4}^{2}\cdot -0.06\approx -125.0235\,\mathrm{N}}\]

\item Aislar posición relativa del centro de presión.

\[\fitmath{\mathrm{xratio}=0.25-\frac{-0.06}{\mathrm{cl}}}\]

\[\fitmath{\mathrm{xratio}=0.25-\frac{-0.06}{(0.6579736267)}\approx 0.3411891\,}\]

\item Convertir posición relativa a distancia.

\[\fitmath{\mathrm{xcp}=\mathrm{xratio}\cdot 1.4}\]

\[\fitmath{\mathrm{xcp}=(0.3411890653)\cdot 1.4\approx 0.4776647\,\mathrm{m}}\]

\item Comprobar escala viscosa.

\[\fitmath{\mathrm{Re}=\frac{1.05\cdot 45\cdot 1.4}{1.75\times10^{-5}}\approx 3780000\,}\]

\end{enumerate}
